\documentclass[13pt,a4paper]{article}
\usepackage[utf8]{inputenc}
\usepackage[T5]{fontenc}
\usepackage[vietnamese]{babel}
\usepackage{geometry}
\geometry{a4paper, margin=1in}
\usepackage{hyperref}
\usepackage{titlesec}
\usepackage{enumitem}
\usepackage{parskip}

\title{\textbf{KỊCH BẢN THUYẾT TRÌNH\\CHUYÊN ĐỀ HÓA HỌC: PHẢN ỨNG CHÁY NỔ}}
\author{\textit{Dành cho Diễn giả}}
\date{}

\begin{document}
\maketitle

\section*{LỜI MỞ ĐẦU}
Xin hân hạnh chào mừng tất cả quý vị, thầy cô và các bạn đã đến với buổi thuyết trình chuyên đề hóa học của chúng ta ngày hôm nay. Hành trình mà chúng ta sắp bước vào không chỉ là những kiến thức khô khan trên sách vở, mà là cả một lăng kính thực tế để giải mã những hiện tượng kì vĩ và nguy hiểm bậc nhất trong tự nhiên: Phản ứng cháy và Nổ.

\hrulefill

\subsection*{SLIDE 1: TIÊU ĐỀ - CHUYÊN ĐỀ HÓA HỌC}
\textbf{Hành động:} (Slide mở màn với hiệu ứng lửa bùng lên mạnh mẽ, tàn lửa bay lơ lửng xung quanh dòng chữ "PHẢN ỨNG CHÁY NỔ").\\
\textbf{Diễn giả:} Chuyên đề ngày hôm nay mang tên: "Phản ứng Cháy Nổ". Chúng ta sẽ cùng nhau khám phá bản chất của sự hủy diệt: Điều kiện hình thành, cơ chế phát triển và những hậu quả thảm khốc không lường trước.

\subsection*{SLIDE 2: KHÁI NIỆM \& ỨNG DỤNG PHẢN ỨNG CHÁY}
\textbf{Hành động:} (Slide hiển thị khối text 3D "1. Phản ứng cháy" cùng thẻ kính chứa nội dung và hình ảnh ngọn lửa uốn lượn).\\
\textbf{Diễn giả:} Xung quanh ta có vô vàn những phản ứng oxi hóa - khử. Vậy phản ứng cháy là gì? Đó chính là những phản ứng oxi hóa - khử mà chất oxi hóa thường là oxygen trong không khí, đặc trưng bởi sự tỏa ra nhiều nhiệt và phát ra ánh sáng. 
Từ xa xưa, tổ tiên đã biết ứng dụng phản ứng này để đốt củi lấy nhiệt, lấy sáng. Và ngày nay, sự cháy của xăng dầu chính là nguồn động lực khổng lồ vận hành động cơ ô tô, máy móc công nghiệp.

\subsection*{SLIDE 3: HẬU QUẢ NGOÀI Ý MUỐN}
\textbf{Hành động:} (Slide Cảnh báo đỏ xuất hiện 2 khối nội dung: "Khí độc ẩn chứa" và "Khí hậu toàn cầu").\\
\textbf{Diễn giả:} Thế nhưng, khi mất kiểm soát, sự cháy biến thành thảm họa. Đầu tiên, nó sinh ra các luồng khí độc lẩn khuất như CO, HCN, HCl, SO2... gây nghẹt thở, ngộ độc trực tiếp cho sinh vật. Thứ hai, và to lớn hơn, lượng khí thải khổng lồ là tác nhân thầm lặng bức tử Trái Đất, gây hiệu ứng nhà kính và biến đổi khí hậu toàn cầu.

\subsection*{SLIDE 4: LÁ PHỔI XANH - RỪNG AMAZON BỐC CHÁY}
\textbf{Hành động:} (Slide với hiệu ứng tàn tro rơi xuống khu rừng Amazon đang đỏ rực).\\
\textbf{Diễn giả:} Không đâu xa, hãy nhìn vào Amazon, lá phổi xanh của hành tinh chúng ta mỗi năm lưu giữ tới 2 tỷ tấn CO2. Đây là nơi sinh sống của 20\% loài thực vật trên Trái Đất. Hàng tỷ cái cây ở đây không chỉ giữ lại lượng carbon khổng lồ mà còn nhả hơi nước tạo mây, điều hòa thời tiết khắp Nam Mỹ. Việc Amazon bốc cháy chính là việc phá vỡ toàn bộ hệ sinh thái mong manh này, giải phóng khí thải nhà kính khổng lồ ngược lại bầu khí quyển.

\subsection*{SLIDE 5: THẢM HỌA THỰC TẾ (VIDEO)}
\textbf{Hành động:} (Giao diện hiển thị viền sáng báo động, phát video "Thảm Họa Thực Tế" - Thảm họa cháy).\\
\textbf{Diễn giả:} Để hình dung rõ hơn, xin mời mọi người hướng lên màn hình để chứng kiến trực tiếp sức tàn phá của ngọn lửa trong thực tế. 
\textit{(Sau video)} Sự hung hãn ấy nuốt chửng mọi thứ trên đường đi của nó!

\subsection*{SLIDE 6: TAM GIÁC CHÁY (TƯƠNG TÁC)}
\textbf{Hành động:} (Hiển thị 3 đỉnh của Tam Giác Cháy. Rọi chuột (hover) vào từng đỉnh để bật chi tiết).\\
\textbf{Diễn giả:} Để một phản ứng cháy bùng phát, bắt buộc phải hội tụ đủ "Tam giác cháy" với 3 yếu tố:
\textit{(Rọi chuột vào điểm đầu tiên)} 1. Phải có Chất cháy: nhiên liệu bắt lửa như xăng, gas, gỗ...
\textit{(Rọi chuột vào điểm thứ hai)} 2. Phải có Oxygen: duy trì sự cháy và nồng độ ưu tiên phải lớn hơn hoặc bằng 14\%.
\textit{(Rọi chuột vào điểm thứ ba)} 3. Phải có Nguồn nhiệt khai mào: một mồi lửa kích hoạt ban đầu.
Chỉ khi hội tụ ĐỒNG THỜI 3 yếu tố này thì ngọn lửa mới tồn tại.

\subsection*{SLIDE 7: CÁC NGƯỠNG NHIỆT ĐỘ KHẮC NGHIỆT}
\textbf{Hành động:} (Chữ NỔ biến mất để lộ ra 3 quả cầu phát sáng. Rọi chuột vào để hiện giải thích).\\
\textbf{Diễn giả:} Và có 3 ngưỡng nhiệt độ đáng chú ý của các hợp chất:
1. \textbf{Điểm chớp nháy}: Nhiệt độ thấp nhất mà hơi chất bốc cháy khi "tiếp xúc" nguồn lửa. Xăng có điểm chớp nháy là -43°C, nên ở nhiệt độ phòng, xăng cực kỳ dễ bắt cháy!
2. \textbf{Nhiệt độ tự bốc cháy}: Là nhiệt độ thấp nhất mà chất tự nổ tung tự cháy mà "KHÔNG CẦN" tiếp xúc trực tiếp nguồn lửa. Với xăng, mốc này là khoảng 247-280°C.
3. \textbf{Nhiệt độ ngọn lửa}: Nhiệt độ cao nhất tạo ra từ vụ cháy. Acetylene cháy tạo ra sức nóng siêu việt 2334°C.

\subsection*{SLIDE 8: NỔ VẬT LÝ}
\textbf{Hành động:} (Giao diện Scan "PHYSICAL ANOMALY DETECTED" và mô hình cái Nồi áp suất hiển thị "PRESSURE CRITICAL 9999 PSI", sau đó vỡ vụn).\\
\textbf{Diễn giả:} Chúng ta thường nói "nổ", nhưng "nổ" chia làm các dạng. Đầu tiên là Nổi vật lý. Đây là sự giãn nở thể tích cực nhanh do tăng áp suất đột ngột mà HOÀN TOÀN KHÔNG kèm theo sự tạo thành chất mới. Nó không hề có phản ứng hóa học. Ví dụ điển hình chính là nổ nồi hơi, nồi áp suất khi chịu tải quá mức.

\subsection*{SLIDE 9: SO SÁNH HÌNH THÁI NỔ}
\textbf{Hành động:} (Giao diện hiện ra 2 thẻ Holographic so sánh Nổ Nồi Áp Suất và Khoang Tàu Dầu).\\
\textbf{Diễn giả:} Các bạn có thể thấy rõ độ tương phản. Nồi áp suất (bên trái) nổ là do áp suất hơi nước tăng vượt giới hạn vỏ nồi - hoàn toàn là Nổ vật lý. Trong khi đó, ở một khoang tàu dầu (bên phải), một tia lửa nhỏ gặp lượng hơi dầu tích tụ sẽ tạo ra phản ứng oxy hóa khử mãnh liệt tạo sức ép to lớn và ngọn lửa tàn phá. Đó gọi là Nổ hóa học.

\subsection*{SLIDE 10: MÔ PHỎNG SỨC ÉP VẬT LÝ (VIDEO)}
\textbf{Hành động:} (Màn hình lưới quét Sci-Fi phát Video Mô Phỏng Sự Nổ).\\
\textbf{Diễn giả:} Mời các bạn xem một đoạn phân tích mô phỏng về sức ép khi giới hạn chịu đựng vật lý bị phá vỡ chớp nhoáng, không hề có sự xuất hiện của lửa. \textit{(Chiếu video)}

\subsection*{SLIDE 11: NỔ HÓA HỌC}
\textbf{Hành động:} (Slide hiển thị mô phỏng hạt nhân đang phân tách - CHEMICAL\_REACTION\_SIMULATION).\\
\textbf{Diễn giả:} Tiếp theo, Nổ hóa học. Định nghĩa rất rõ ràng: Sự nổ gây ra do phản ứng hóa học diễn ra với tốc độ cực nhanh, tỏa rất nhiều nhiệt, sinh ra lượng khí lớn làm tăng áp suất đột ngột. Điều kiện BẮT BUỘC của nó là phải có sự sinh ra chất mới do biến đổi hóa học.

\subsection*{SLIDE 12: VÍ DỤ THỰC TẾ - KHOANG TÀU DẦU (TƯƠNG TÁC)}
\textbf{Hành động:} (Rê chuột thành tiêu điểm quét nhiệt (Thermal Scan) để lộ đám khí dầu đỏ rực ẩn dưới hầm tàu tối tăm).\\
\textbf{Diễn giả:} \textit{(Rê chuột)} Các bạn hãy nhìn hình ảnh quét nhiệt này. Dù khoang tàu đã hút cạn dầu bẩn, một lượng hơi bám dính vẫn tích tụ thành mây khí ở dưới đáy. Lượng hơi này trộn với Oxygen không khí thành "hỗn hợp dễ nổ". Một mỏ hàn sửa chữa có thể lập tức tiễn mọi thứ tan tành mây khói.

\subsection*{SLIDE 13: PHẢN ỨNG METHANE}
\textbf{Hành động:} (Phân tử CH4 và 2O2 va chạm nổ lóa mắt tạo thành CO2 và 2H2O).\\
\textbf{Diễn giả:} Khí Methane trong vùng kín bắt lửa là một thảm họa khu mỏ. Phương trình: $CH_4 + 2O_2 \rightarrow CO_2 + 2H_2O$. Sự đứt lìa của các kết nối hóa học cũ để sinh ra chất mới giải phóng năng lượng siêu cường, $\Delta H_{298}^\circ = -802.5 \text{ kJ}$, tạo ra hiệu ứng hủy diệt kinh hoàng.

\subsection*{SLIDE 14: SỰ NỔ CỦA TNT (TƯƠNG TÁC)}
\textbf{Hành động:} (Màn hình rung lắc cảnh báo Class 1 Explosive. Click vào nút đỏ DETONATE giữa màn hình để nổ bung ra phương trình).\\
\textbf{Diễn giả:} Khi nhắc đến nổ hóa học, không ai không biết đến công thức thuốc nổ huyển thoại TNT (Trinitrotoluene). \textit{(Click nổ)} Bùmmmm! Sự phân hủy tức thời sinh ra CO, H2, N2 đẩy thể tích khí giãn nở khủng khiếp, phá vỡ hoàn toàn mọi thứ cản đường.

\subsection*{SLIDE 15: HỖN HỢP NỔ HYDROGEN \& OXYGEN (TƯƠNG TÁC)}
\textbf{Hành động:} (Mô phỏng lò phản ứng chứa siêu vi hạt H2 O2 bay vòng quanh. Click mạnh vào lõi lò để kích nổ).\\
\textbf{Diễn giả:} Nhưng đứng đầu về sức nổ là gì? Xin mời nhìn vào buồng chứa hỗn hợp H2 và O2 này. \textit{(Click)} Phản ứng $2H_2 + O_2 \rightarrow 2H_2O$. Hydrogen và Oxi với thể tích vàng 2:1 là hỗn hợp nổ mạnh nhất! Nhiệt lượng giải tỏa khổng lồ cỡ $-571.6 \text{ kJ}$ khiến nó tuyệt vời đến mức được ứng dụng làm nhiên liệu cho tên lửa du hành không gian.

\subsection*{SLIDE 16: CẢNH BÁO BÓNG BAY HYDROGEN}
\textbf{Hành động:} (Các quả bóng bay chứa dòng cảnh báo bay lên từ dưới màn hình).\\
\textbf{Diễn giả:} Nhưng chính loại khí nhẹ, nổ mạnh này lại thường bị con người bất cẩn dùng để bơm... Bóng bay cho trẻ con. Khí Hydro nếu gặp môi trường ngoài trời nóng rực hoặc một tàn thuốc, vụ nổ hóa học sẽ khiến người cầm bóng bỏng nặng. Do đó, khuyến cáo bắt buộc là dùng khí hiếm an toàn Helium (He) cho các quả bóng bay trang trí!

\subsection*{SLIDE 17: NỔ BỤI (TƯƠNG TÁC KÍCH NỔ)}
\textbf{Hành động:} (Mô phỏng hàng ngàn hạt bụi bay lơ lửng. Rê chuột làm bụi dạt ra. Click kích mồi lửa tạo quả cầu lửa dây chuyền).\\
\textbf{Diễn giả:} Và dạng nổ cuối cùng, ám ảnh nhất bởi sự bất ngờ của nó ở kho xưởng: Nổ bụi. Nó là khi bụi bặm vô hại lơ lửng trong không khí đủ đặc tạo thành "mây bụi". \textit{(Click bụi)} Chỉ một cú nổ kích mồi, ngọn lửa lan truyền dây chuyền liên tục từ hạt bụi này sang hạt bụi khác xuyên qua không khí với tốc độ cấp số nhân, tạo thành một cầu lửa khổng lồ!

\subsection*{SLIDE 18: 5 YẾU TỐ CỦA NỔ BỤI}
\textbf{Hành động:} (Slide mở ra 5 thẻ bài 3D bay lên tương ứng với 5 yếu tố nổ bụi).\\
\textbf{Diễn giả:} Và để nổ bụi xảy ra, nó khắc nghiệt hơn cháy ở chỗ phải thỏa mãn tận 5 yếu tố: (1) Hạt phải là Bụi Mịn, (2) Bụi lơ lửng với Nồng Độ Cao, (3) Có Oxygen, (4) Có một Nguồn Nhiệt khởi mào, và (5) Phải ở trong Không Gian Kín (để tạo sức ép). Ghi nhớ kỹ: Thiếu MỘT trong năm yếu tố này, NỔ BỤI KHÔNG bao giờ xảy ra!

\subsection*{SLIDE 19: ĐIỂM DANH "SÁT THỦ" BỤI}
\textbf{Hành động:} (Bảng điểm danh 4 ngành công nghiệp nổ bụi. Rê chuột vào từng ngành để soi sáng thẻ bài tương ứng).\\
\textbf{Diễn giả:} Đừng tưởng chỉ thuốc súng mới nổ. Những thứ giản dị quanh ta mới là sát thủ! Trong ngành công nghiệp Thực Phẩm là bột mì, đường, bột sữa. Ngành Chế biến là mùn cưa. Đặc biệt nguy hiểm nhất là Gia công Kim Loại (bột nhôm, bột magie) – những loại này khi nổ sẽ tạo ra nhiệt lượng và áp suất hủy hoại mức độ tối đa!

\subsection*{SLIDE 20: TỔNG KẾT \& TRÒ CHƠI}
\textbf{Hành động:} (Hệ thống báo màn hình "Nhiệm vụ hoàn tất". Click nút "KHÓA HUẤN LUYỆN THỰC TẾ" để loading vào Game).\\
\textbf{Diễn giả:} Vâng xin cảm ơn quý thính giả, chúng ta đã đi trọn vẹn từ phản ứng cháy, nổ vật lí, nổ hóa học đến nổ bụi. Các bạn đã nắm trọn bản chất sự hủy diệt trong tay. 
Và đây là lúc kiểm chứng kiến thức đó! Mời các bạn tham gia "Khóa Huấn Luyện Thực Tế" - trò chơi tương tác dập lửa giải cứu phòng lab ảo ngay bây giờ!
\textit{(Click nút Play - Game bắt đầu).}

\end{document}
